\section{策略梯度与 Actor-Critic 结合}
\subsection{Actor-Critic 架构导入}
Actor-Critic 结构将策略（actor）与 value function（critic）一起训练。Critic 提供 advantage 估计，actor 则根据 advantage 调整策略，形成双塔结构：
\begin{enumerate}
    \item Actor 产出 action distribution $\pi_\theta(a\mid s)$。
    \item Critic 估计 $V_w(s)$，供 advantage 计算 $\hat{A}_t = G_t - V_w(s_t)$。
    \item Actor 更新使用 above advantage，Critic 用 TD loss 更新值网络。
\end{enumerate}
\begin{knowledgebox}{Actor-Critic 的协同}
Actor 根策略提供 high-level decision；Critic 判断 reward 估值，两者共同缓解传统策略梯度的高 variance。
\end{knowledgebox}

Critic 的 bootstrap update 用 $\delta = r_t + \gamma V_w(s_{t+1}) - V_w(s_t)$ 定义 TD error，再最小化 $\delta^2$ 能让 critic 提供更准确的 baseline。

\subsection{Actor-Critic vs 原始策略梯度}
原始策略梯度仅使用 Monte Carlo reward，因此 variance 大、数据利用差；Actor-Critic 则引入 value estimate 作为 baseline，降低 variance 并实现 bootstrapping。
\begin{table}[H]
\centering
\small
\begin{tabular}{p{0.34\textwidth}p{0.34\textwidth}p{0.3\textwidth}}
\toprule
方法 & 估计类型 & 样本利用 \\
\midrule
REINFORCE & Fully Monte Carlo & 每条轨迹使用一次 \\
Actor-Critic & Bootstrapped value & 同一数据可反复更新 \\
PPO & Clipped surrogate + value head & 利用多个 epoch \\
\bottomrule
\end{tabular}
\caption{策略梯度系列方法对照}
\end{table}

在实际训练中，actor 的 gradient 表达式依旧是 $\mathbb{E}[\nabla_\theta \log \pi_\theta(a\mid s) \cdot \hat{A}]$，但 advantage 由 critic 提供，可用 GAEn 形式提高稳定。Critic 更新时，可通过多步 TD（如 TD($\lambda$)）融合 Monte Carlo 和 bootstrap 信号，在双塔架构中形成协同。

\subsection{本章小结}
Actor-Critic 把策略梯度与 value function 结合，降低了 variance 并提升了样本效率；PPO/TRPO 进一步在此基础上加入 trust region，形成更稳定的 training loop。

